% =====================================================================
%  PREÁMBULO: paquetes y configuración
% =====================================================================

% ---------- Idioma y codificación ----------
\usepackage[utf8]{inputenc}
\usepackage[T1]{fontenc}
\usepackage[english, catalan, spanish, es-noquoting, es-tabla]{babel}   % el último idioma es el principal; es-noquoting: "<<" no abre comillas, las comillas latinas se escriben « » directamente; es-tabla: «Tabla» y no «Cuadro» (guía EPS 4.2)
\usepackage{csquotes}
\addto\captionsspanish{\renewcommand{\contentsname}{Índice de contenidos}}   % nombre de la guía EPS 2.5

% ---------- Márgenes ----------
\usepackage[a4paper,top=2.5cm,bottom=2.5cm,left=3cm,right=3cm,inner=3cm,outer=2.5cm,marginparwidth=2cm]{geometry}

% ---------- Interlineado ----------
% Interlineado 1,5 como el de Word en TODO el documento: texto, notas, pies de figura, tablas y bibliografía.
% No se usa setspace porque este deja en sencillo las notas y los flotantes. LaTeX ya separa las líneas
% 1,2 veces el tamaño de letra, así que 1,24 x 1,2 ≈ 1,5 líneas.
\linespread{1.24}

% ---------- Matemáticas y listas ----------
\usepackage{amsmath}
\usepackage{enumitem}

% ---------- Tipos de letra (guía EPS 1.2) ----------
% Familia 2 (Palatino) para el texto, los títulos y las fórmulas (en itálica); familia 1 (TeX Gyre Heros,
% clon libre de Helvética) para la bibliografía; familia 3 (Courier) para el código.
% La guía solo «recomienda» la familia 1 para los títulos; se dejan en Palatino, igual que el texto.
\usepackage{newpxtext}                   % también pone TeX Gyre Heros como familia sin serifa
\usepackage{newpxmath}
\renewcommand{\ttdefault}{pcr}

% ---------- Formato de títulos ----------
\usepackage{titlesec}
\titleformat{\paragraph}[block]{\normalfont\large\bfseries}{\theparagraph}{1em}{}

% ---------- Figuras ----------
\usepackage{graphicx}
\usepackage{float}
\usepackage{caption}                     % \captionsetup para ajustar los pies de una figura concreta
% Colocación de figuras: que compartan página con el texto en lugar de quedarse solas en una
% página con huecos arriba y abajo
\renewcommand{\topfraction}{0.85}        % hasta el 85 % de una página puede ser figuras arriba
\renewcommand{\bottomfraction}{0.6}
\renewcommand{\textfraction}{0.1}        % basta un 10 % de texto para que la página no sea solo de figuras
\renewcommand{\floatpagefraction}{0.85}  % página solo de figuras únicamente si estas la llenan casi entera
\setcounter{topnumber}{3}
\setcounter{totalnumber}{4}

% Ningún párrafo se parte dejando menos de tres líneas al final o al principio de una página
\usepackage[defaultlines=3, all]{nowidow}
\raggedbottom                            % a dos caras LaTeX estira el espacio vertical para igualar el pie de las páginas; con figuras [H] deja huecos
\usepackage{eso-pic}                     % imagen de fondo en la portada
\usepackage[absolute, overlay]{textpos}

% ---------- Gráficas (pgfplots) ----------
\usepackage{pgfplots}
\pgfplotsset{compat=1.18}
\usetikzlibrary{positioning, arrows.meta, fit, backgrounds, shapes.geometric}   % esquemas de bloques
\definecolor{serie1}{HTML}{2A78D6}       % azul: serie principal
\definecolor{serie2}{HTML}{EB6834}       % naranja: segunda serie
\definecolor{tintasecundaria}{HTML}{52514E}
\pgfplotsset{
	grafica/.style={
		width=0.92\linewidth, height=6cm,
		axis lines*=left, axis line style={tintasecundaria},
		tick style={tintasecundaria}, tick label style={font=\small, text=tintasecundaria},
		label style={font=\small}, grid style={gray!20}, ymajorgrids,
		/pgf/number format/use comma, /pgf/number format/1000 sep={\,},
		legend style={draw=none, font=\small, legend columns=-1},
	},
}

% ---------- Tablas ----------
\usepackage{longtable}
\usepackage{tabularx}
\usepackage{booktabs}
\usepackage{eurosym}                     % \euro en las tablas de precios

% ---------- Páginas en horizontal ----------
\usepackage{pdflscape}                   % \begin{landscape} ... \end{landscape}
\usepackage{rotating}
\usepackage{afterpage}                   % \afterpage{...}: la página horizontal sale tras la actual, sin dejar hueco

% ---------- Cajas, colores y notas ----------
\usepackage{xcolor}
\usepackage[most]{tcolorbox}
\usepackage{todonotes}                   % \todo{...} para notas pendientes

% Marca visible para las citas que aún faltan; buscar "\citapendiente" antes de entregar
\newcommand{\citapendiente}{\unskip\textsuperscript{\textcolor{red}{[cita]}}}

% Nota de fuente al pie de una tabla o figura (guía EPS 4.3). Uso: \fuente{elaboración propia}
\newcommand{\fuente}[1]{\par\smallskip{\footnotesize Fuente: #1.}\par}

% ---------- Código fuente ----------
\usepackage{listings}
\renewcommand{\lstlistingname}{Código}
\renewcommand{\lstlistlistingname}{Índice de código}

\lstdefinelanguage{JavaScript}{
	keywords={const, let, var, function, return, async, await, export, import, from, for, of, if, else, new, class},
	keywordstyle=\color{blue}\bfseries,
	comment=[l]{//},
	morecomment=[s]{/*}{*/},
	commentstyle=\color{gray}\itshape,
	string=[b]',
	morestring=[b]",
	stringstyle=\color{teal},
	sensitive=true
}

\lstset{
	basicstyle=\ttfamily\footnotesize,
	breaklines=true,
	breakindent=0pt,
	columns=fullflexible,
	keepspaces=true,
	captionpos=b,
	showstringspaces=false,
	literate={á}{{\'a}}1 {é}{{\'e}}1 {í}{{\'i}}1 {ó}{{\'o}}1 {ú}{{\'u}}1 {ñ}{{\~n}}1
	         {Á}{{\'A}}1 {É}{{\'E}}1 {Í}{{\'I}}1 {Ó}{{\'O}}1 {Ú}{{\'U}}1 {Ñ}{{\~N}}1
}

\lstdefinestyle{codigo}{
	frame=single,
	numbers=left,
	numberstyle=\tiny\color{gray},
	keywordstyle=\color{blue},
	stringstyle=\color{teal},
	commentstyle=\color{gray},
	tabsize=4
}

% ---------- Enlaces ----------
\usepackage[hidelinks, hypertexnames=false]{hyperref}   % la portada y los preliminares repiten números de página

% ---------- Referencias, bibliografía ----------
\usepackage[backend=biber, style=authoryear, maxcitenames=2, sorting=nyt]{biblatex}
\addbibresource{referencias.bib}
% Las URL largas de la bibliografía pueden partirse en cualquier carácter (si no, se salen del margen)
\setcounter{biburlnumpenalty}{9000}
\setcounter{biburlucpenalty}{9000}
\setcounter{biburllcpenalty}{9000}
\appto{\bibsetup}{\sloppy\raggedright}   % sin justificar: las URL que no se pueden partir dejaban líneas muy estiradas
\renewcommand*{\bibfont}{\sffamily}      % guía EPS 2.8: la bibliografía va en familia 1

% ---------- Glosario (después de hyperref) ----------
\PassOptionsToPackage{lang-warn=false}{datatool-base}   % datatool no trae módulos de idioma instalados
\usepackage[toc, acronym, nolangwarn]{glossaries}   % nolangwarn: no hay módulo de glossaries para catalán
\makeglossaries

% ---------- Encabezado y pie ----------
\usepackage{fancyhdr}
\setlength{\headheight}{14pt}            % el cuerpo va en \large y la cabecera no cabía en 12pt
\pagestyle{fancy}
\fancyhf{}
% Guía EPS 1.4: a dos caras, con cabeceras distintas en páginas pares e impares.
% El número va siempre en el borde exterior; el capítulo en las pares y la sección en las impares.
\fancyhead[LE,RO]{\thepage}
\fancyhead[RE]{\nouppercase{\leftmark}}
\fancyhead[LO]{\nouppercase{\rightmark}}
\renewcommand{\headrulewidth}{0.4pt}
\usepackage{emptypage}                   % las páginas en blanco que añade el modo a dos caras, sin cabecera

% Estilo para las páginas que abren capítulo
\fancypagestyle{plain}{
	\fancyhf{}
	\fancyfoot[C]{\thepage}
	\renewcommand{\headrulewidth}{0pt}
}
